% Part: first-order-logic
% Chapter: tableaux
% Section: quantifier-rules

\documentclass[../../../include/open-logic-section]{subfiles}

\begin{document}

\olfileid{fol}{tab}{qrl}

\olsection{संख्यापकांचे नियम}

\subsection{$\lforall$ साठीचे नियम}

\begin{defish}
\AxiomC{\sFmla{\True}{\lforall[x][!A(x)]}}
\RightLabel{\TRule{\True}{\forall}}
\UnaryInfC{\sFmla{\True}{!A(t)}}
\DisplayProof
\hfill
\AxiomC{\sFmla{\False}{\lforall[x][!A(x)]}}
\RightLabel{\TRule{\False}{\lforall}}
\UnaryInfC{\sFmla{\False}{!A(a)}}
\DisplayProof
\end{defish}

\TRule{\True}{\lforall} मध्ये, $t$ हे बंद पद आहे (म्हणजे, चरे नसलेले
पद). \TRule{\False}{\lforall} मध्ये, $a$~हा असा !!a{constant} आहे की तो
\TRule{\False}{\lforall} नियमाच्या वरच्या शाखेत कोठेही येता कामा नये.
\TRule{\False}{\forall} अनुमानाचा $a$ हा \emph{आयगेन चर} आहे, असे आपण
म्हणतो.\footnote{वरील नियमात $a$ हा !!a{constant} असला तरी आपण
``आयगेन चर'' ही संज्ञा वापरतो. यामागे ऐतिहासिक कारणे आहेत.}

\subsection{$\lexists$ साठीचे नियम}

\begin{defish}
\AxiomC{\sFmla{\True}{\lexists[x][!A(x)]}}
\RightLabel{\TRule{\True}{\lexists}}
\UnaryInfC{\sFmla{\True}{!A(a)}}
\DisplayProof
\hfill
\AxiomC{\sFmla{\False}{\lexists[x][!A(x)]}}
\RightLabel{\TRule{\False}{\lexists}}
\UnaryInfC{\sFmla{\False}{!A(t)}}
\DisplayProof
\end{defish}

पुन्हा, $t$~हे बंद पद आहे आणि $a$~हा असा !!a{constant} आहे की तो
\TRule{\True}{\lexists} नियमाच्या वरच्या शाखेत येत नाही. आपण $a$ ला
\TRule{\True}{\lexists} अनुमानाचा \emph{आयगेन चर} म्हणतो.

\TRule{\False}{\lforall} किंवा \TRule{\True}{\lexists} अनुमानाच्या वरच्या
शाखेत आयगेन चर येत नसावा, या अटीला \emph{आयगेन चराची अट} म्हणतात.

\begin{explain}
ही प्रथा आठवा: $!A$ हे असे !!a{formula} असेल की त्यात
!!{variable}~$x$ मुक्त असेल, तर हे आपण~$!A(x)$ असे लिहून दर्शवतो.
त्याच संदर्भात, मग $!A(t)$ हा~$\Subst{!A}{t}{x}$ चा संक्षेप असतो.
त्यामुळे \TRule{\False}{\lexists} नियम पुढीलप्रमाणेही लिहिता येईल:
\begin{prooftree}
  \AxiomC{\sFmla{\False}{\lexists[x][!A]}}
  \RightLabel{\TRule{\False}{\lexists}}
  \UnaryInfC{\sFmla{\False}{\Subst{!A}{t}{x}}}
\end{prooftree}
लक्षात घ्या की~$!A$ मध्ये $t$ आधीच येऊ शकते; उदा., $!A$ हे
~$\Atom{\Obj P}{t,x}$ असू शकते. म्हणून,
$\sFmla{\False}{\lexists[x][\Atom{\Obj P}{t,x}]}$ पासून
$\sFmla{\False}{\Atom{\Obj P}{t,t}}$ निष्पन्न करणे हा
\TRule{\False}{\lexists} चा योग्य उपयोग आहे. परंतु
\TRule{\False}{\lforall} आणि~\TRule{\True}{\lexists} मधील आयगेन चराच्या
अटींनुसार !!{constant}~$a$ हा~$!A$ मध्ये येता कामा नये. म्हणून,
$\sFmla{\False}{\lforall[x][\Atom{\Obj P}{a,x}]}$ पासून
$\sFmla{\False}{\Atom{\Obj P}{a,a}}$ हे
~$\TRule{\False}{\lforall}$ वापरून योग्य रीतीने निष्पन्न करता येत नाही.
\end{explain}

\begin{explain}
\TRule{\True}{\lforall} आणि \TRule{\False}{\lexists} मध्ये पद~$t$ बंद
असण्याव्यतिरिक्त त्यावर आणखी कोणतेही निर्बंध नाहीत. याउलट,
\TRule{\True}{\lexists} आणि \TRule{\False}{\lforall} नियमांमध्ये आयगेन
चराच्या अटीनुसार संबंधित अनुमानांच्या वरच्या शाखांमध्ये !!{constant}~$a$
कोठेही येता कामा नये. ही पद्धत निर्दोष राहील याची खात्री करण्यासाठी ही अट
आवश्यक आहे. ही अट नसती, तर पुढील आकृती
$\lexists[x][\formula{A}(x)] \lif \lforall[x][\formula{A}(x)]$ साठीचा
बंद !!{tableau} ठरली असती:
\begin{center}
\begin{tableau}{}
  [\sFmla{\False}{\lexists[x][\formula{A}(x)] \lif \lforall[x][\formula{A}(x)]}, just=\TAss
    [\sFmla{\True}{\lexists[x][\formula{A}(x)]},
      just={\TRule{\False}{\lif}[1]}
      [\sFmla{\False}{\lforall[x][\formula{A}(x)]},
        just={\TRule{\False}{\lif}[1]}
        [\sFmla{\True}{\formula{A}(a)},
          just={\TRule{\True}{\lexists}[2]}
          [\sFmla{\False}{\formula{A}(a)},
            just={\TRule{\False}{\lforall}[3]}, close]
        ]
      ]
    ]
  ]
\end{tableau}
\end{center}
परंतु $\lexists[x][\formula{A}(x)] \lif
\lforall[x][\formula{A}(x)]$ हे वैध नाही.
\end{explain}

\end{document}
